\section{Scheme Design}
\label{sec:scheme-design}

\subsection{Sample Preparation and Matrix}
The test matrix was prepared as a single 400~L bulk batch of deionised water reconstituted with a
certified hardness/alkalinity concentrate to approximate a distribution-network sample, then
gravimetrically spiked with traceable single-element standards for Pb, Cd, Cu, and Cr. The batch
was homogenised by recirculation for 4~hours before bottling into 120~mL HDPE bottles, with two
bottles (A, B) allocated per participant.

\subsection{Homogeneity and Stability Testing}
Ten bottles were selected at random from the filling run and analysed in duplicate by Meridian's
reference laboratory (ICP-MS, validated to ISO 17294-2). The between-bottle standard deviation
(\(s_{bb}\)) for each analyte was compared against the fitness-for-purpose target per ISO 13528
Annex B:

\begin{center}
\small
\begin{tabular}{@{}lccc@{}}
\toprule
\textbf{Analyte} & \textbf{Reference mean (\(\mu\)g/L)} & \(s_{bb}\) \textbf{(\(\mu\)g/L)} & \textbf{Homogeneity verdict} \\
\midrule
Lead (Pb)     & 8.19 & 0.071 & \perfsat \\
Cadmium (Cd)  & 2.04 & 0.028 & \perfsat \\
Copper (Cu)   & 145.6 & 1.42 & \perfsat \\
Chromium (Cr) & 12.7 & 0.118 & \perfsat \\
\bottomrule
\end{tabular}
\end{center}

Stability was confirmed by re-analysing retained bottles after 21~days of storage at the
temperature profile experienced by participants in transit; no analyte drifted by more than 0.3 of
\(\sigma_{pt}\), satisfying the stability criterion of ISO 13528 \S B.3.

\subsection{Assigned Value and Target Standard Deviation}
Assigned values (\(x_{pt}\)) were computed as the robust mean of all 14 participant results per
analyte using Algorithm A (ISO 13528 Annex C), chosen over a formal reference-value approach
because reference-laboratory data were used only for homogeneity/stability confirmation and not
disclosed prior to the round. The proficiency standard deviation (\(\sigma_{pt}\)) was fixed a
priori as a fitness-for-purpose percentage of the assigned value, informed by the Horwitz-type
concentration dependence of inter-laboratory precision \cite{horwitz1980method} rather than
derived from the observed participant spread.

\begin{center}
\small
\begin{tabular}{@{}lccc@{}}
\toprule
\textbf{Analyte} & \(x_{pt}\) \textbf{(\(\mu\)g/L)} & \(\sigma_{pt}\) \textbf{(\(\mu\)g/L)} & \textbf{Basis for \(\sigma_{pt}\)} \\
\midrule
Lead (Pb)     & 8.20  & 0.65 & 8\% fitness-for-purpose \\
Cadmium (Cd)  & 2.05  & 0.16 & 8\% fitness-for-purpose \\
Copper (Cu)   & 146.0 & 14.6 & 10\% fitness-for-purpose \\
Chromium (Cr) & 12.8  & 1.02 & 8\% fitness-for-purpose \\
\bottomrule
\end{tabular}
\end{center}

\subsection{Round Timeline}
\begin{center}
\small
\begin{tabular}{@{}ll@{}}
\toprule
\textbf{Milestone} & \textbf{Date} \\
\midrule
Bulk batch preparation and homogeneity testing & 24 February 2026 \\
Sample dispatch to participants & 10 March 2026 \\
Results submission deadline & 07 April 2026 \\
Statistical evaluation and Algorithm A computation & 09--14 April 2026 \\
Draft report technical review & 16 April 2026 \\
Individual reports issued to participants & 21 April 2026 \\
\bottomrule
\end{tabular}
\end{center}
